\documentclass[10pt,a4paper]{amsart}
\usepackage[margin=19mm]{geometry}
\usepackage{amsmath,amssymb,mathtools}
\usepackage[scr=rsfs,cal=euler]{mathalfa}
\usepackage{xfrac}
\usepackage[unicode,hidelinks,bookmarks=false]{hyperref}
\newcommand{\ie}{i.e.\ }
\newcommand{\cf}{cf.\ }
\newcommand{\resp}{respectively\ }
\newcommand{\Subsection}{\subsection}
\providecommand{\nobreakdash}{\nobreak\hbox{-}}
\setlength{\emergencystretch}{2em}
\multlinegap0pt
\usepackage{amssymb}
\usepackage{float}
\newtheorem{theorem}{Theorem}
\newtheorem{lemma}{Lemma}
\title{A New Classification of Positive Integers Via New Divisor Functions}
\author{Brahim Mittou}
\date{Source: arXiv:2509.08844v3 (CC BY 4.0)}
\begin{document}
\maketitle

\noindent\emph{Source and licence.} A New Classification of Positive Integers Via New Divisor Functions. Version arXiv:2509.08844v3 (CC BY 4.0), distributed by arXiv under the Creative Commons Attribution 4.0 International licence, http://creativecommons.org/licenses/by/4.0/, as recorded in the arXiv licence metadata for this exact version and shown on the abstract page https://arxiv.org/abs/2509.08844v3. Full licence notice: \texttt{licenses/arxiv-2509.08844-CC-BY-4.0.txt}.\par\smallskip

\noindent\textit{Editorial scope.} Counted unit: the article's theorem thm:factorization (source0.tex lines 572--601). The article's complete original proof of it is delivered with it (source0.tex lines 778--882, one \texttt{proof} environment of 105 lines, which the article places away from the statement). No mathematics is written, repaired or re-derived: the statement and the proof are the article's own lines, in the article's own notation, and no retained line is substituted or re-flowed.  Retained uncounted as the source-local context the proof needs: lines 604--621; lines 721--770; lines 732--738; lines 764--769.  Nothing was omitted from any retained span, so the package carries no omission record.  No retained line contains a citation, so the wrapper carries no bibliography and no bibliography entry.  No retained line contains a figure environment, an image-inclusion command or drawing code, so no image is delivered and the package declares no asset dependency.  Editorial formatting only, in addition: an AMS \texttt{amsart} preamble that restates the article's own declarations and macros used by the retained lines, namely the environments \texttt{theorem}, \texttt{lemma} on separate counters, together with the standard packages those lines need.  The retained environments print in this document as Theorem 1, Lemma 1 in order of appearance, whereas the article prints the same items with the numbering of its own full document.  The complete native source of the article as distributed in the arXiv e-print for this exact version is bundled unchanged as \texttt{sources/184852/source0.tex}.
\begin{lemma}\label{lem:determinant}
Let $f(x)$ and $g(x)$ be the polynomials:
\[
f(x)=a_2x^2+a_1x+a_0,\quad
g(x)=b_2x^2+b_1x+b_0,
\]
where the coefficients belong to any commutative ring. Then
\[
f(x)g(y)-f(y)g(x)
=
(x-y)
\left(
C_2xy+C_1(x+y)+C_0
\right),
\]
where
$C_2=a_2b_1-a_1b_2$, $C_1=a_2b_0-a_0b_2$, and $C_0=a_1b_0-a_0b_1$.
\end{lemma}

\begin{lemma}\label{lem:closedcoeff}
Let \(m\ge2\) be an integer, and let \(p,\ q\) be primes satisfying $p<q<p^2$ and $q^2<p^3$.
Write
\[
\sigma_o(p^{2m}q^2)=A_2^{(m)}q^2+A_1^{(m)}q+A_0^{(m)},
\]
and
\[
\sigma_e(p^{2m}q^2)=B_2^{(m)}q^2+B_1^{(m)}q+B_0^{(m)}.
\]
Then
\begin{equation}\label{E8}
A_2^{(m)}
=
p^{2m}
+
\frac{p^{2m-1}-p}{p^2-1},
\end{equation}
\begin{equation}\label{E9}
A_1^{(m)}
=
1+p^{2m}
+
\frac{p^{2m+1}-p}{p^2-1},
\end{equation}
\begin{equation}\label{E10}
A_0^{(m)}
=
1+
\frac{p^{2m+1}-p^3}{p^2-1},
\end{equation}
\begin{equation}\label{E11}
B_2^{(m)}
=
p^{2m-1}
+
\frac{p^{2m}-1}{p^2-1},
\end{equation}
\begin{equation}\label{E12}
B_1^{(m)}
=
\frac{p^{2m}-p^2}{p^2-1},
\end{equation}
\begin{equation}\label{E13}
B_0^{(m)}
=
p+
\frac{p^{2m+2}-p^2}{p^2-1}.
\end{equation}
\end{lemma}

\begin{equation}
A_2^{(m)}
=
p^{2m}
+
\frac{p^{2m-1}-p}{p^2-1},
\end{equation}

\begin{equation}
B_0^{(m)}
=
p+
\frac{p^{2m+2}-p^2}{p^2-1}.
\end{equation}

\begin{theorem}\label{thm:factorization}
Let \(m\) be a positive integer and let \(p<q<r\) be primes satisfying $r<p^2$ and $r^2<p^3$.
Define
\[
N_{2m}
=
\sigma_e(p^{2m}q^2)\sigma_o(p^{2m}r^2)
-
\sigma_e(p^{2m}r^2)\sigma_o(p^{2m}q^2).
\]
Then
\[
N_{2m}
=
(r-q)\,
A(p,q,r)\,
F_m(p),
\]
where
\[
A(p,q,r)=qr+q+r-p(q+r+1),
\]
and
\[
F_m(p)
=
\frac{p^{4m}-1}{p-1}+p^{2m-1}(p+1).
\]
In particular, if $A(p,q,r)=0$, then $k(p^{2m}q^2)=k(p^{2m}r^2)$.
\end{theorem}

\begin{proof}[Proof of Theorem \ref{thm:factorization}]
For \(m=1\), the positive divisors of \(p^2q^2\) and \(p^2r^2\) are
\[
1,\; p,\; q,\; p^2,\; pq,\; q^2,\; p^2q,\; pq^2,\; p^2q^2,
\]
and
\[
1,\; p,\; r,\; p^2,\; pr,\; r^2,\; p^2r,\; pr^2,\; p^2r^2,
\]
respectively. A direct computation gives
\[
\begin{aligned}
&\sigma_e(p^{2}q^2)\sigma_o(p^{2}r^2)
-\sigma_e(p^{2}r^2)\sigma_o(p^{2}q^2) \\
&\qquad
=(r-q)\bigl(-pq-pr+qr-p+q+r\bigr)
\left(p^3+2p^2+2p+1\right) \\
&\qquad
=(r-q)\,A(p,q,r)\,F_1(p).
\end{aligned}
\]
Hence the asserted factorization holds when \(m=1\).
Now assume that \(m\ge2\). By Lemmas \ref{lem:determinant} and \ref{lem:closedcoeff},
\begin{equation}\label{E14}
N_{2m}
=
(q-r)
\left(
C_2^{(m)}qr
+
C_1^{(m)}(q+r)
+
C_0^{(m)}
\right),
\end{equation}
where
\[
C_2^{(m)}=A_2^{(m)}B_1^{(m)}-A_1^{(m)}B_2^{(m)},
\]
\[
C_1^{(m)}=A_2^{(m)}B_0^{(m)}-A_0^{(m)}B_2^{(m)},
\]
and
\[
C_0^{(m)}=A_1^{(m)}B_0^{(m)}-A_0^{(m)}B_1^{(m)}.
\]
Using Identities \eqref{E8}--\eqref{E13}, we obtain
\[
C_2^{(m)}
=
-
\left(
\frac{p^{4m}-1}{p-1}
+
p^{2m-1}(p+1)
\right)
=-F_m(p),
\]
\[
C_1^{(m)}
=
p^{4m}+p^{2m-1}(p^2-1)-1
=-(p-1)C_2^{(m)},
\]
and
\[
C_0^{(m)}
=
p
\left(
\frac{p^{4m}-1}{p-1}
+
p^{2m-1}(p+1)
\right)
=-pC_2^{(m)}.
\]
Substituting these identities into \eqref{E14} gives
\[
\begin{aligned}
N_{2m}
&=(q-r)
\left(
C_2^{(m)}qr
+C_1^{(m)}(q+r)
+C_0^{(m)}
\right) \\
&=(q-r)C_2^{(m)}
\left(
qr-(p-1)(q+r)-p
\right) \\
&=(r-q)\,
A(p,q,r)\,
F_m(p),
\end{aligned}
\]
since
\[
qr-(p-1)(q+r)-p
=
qr+q+r-p(q+r+1)
=
A(p,q,r).
\]
This completes the proof.
\end{proof}

\end{document}
